\section{自然语言处理}
\subsection{单词表示}
自然语言处理（Natural Language Processing，NLP），它的目标就是让计算机理解人说的话，进而帮助人们更好地完成某些事情。
语言是由文字构成的，而语言的含义是由单词构成的，即单词是含义的最小单位。所以，为了让计算机理解自然语言，第一步就是让它理解单词含义。

\subsection{基于同义词词典的方法}
要表示单词含义，首先可以考虑通过人工方式来定义单词含义，一种方式就是像词典一样，一个词一个词地说明单词含义。
NLP中广泛运用的一种词典是同义词词典（thesaurus），在词典中，含义类似的单词被归类到同一个组中。
比如， car 的同义词有automobile、motorcar等。另外，在同义词词典中，有时还会定义更多的关系，比如上位下位关系、整体部分关系等。
在NLP领域，最著名的同义词词典是普林斯顿大学于 1985 年开始开发的同义词词典 WordNet。

同义词词典有许多缺陷，第一点是它难以实时更新，随着时间的推移，会有新词不断出现，也会有旧词被赋予新的含义，
同义词词典很难反映出这种变化。第二点是生成同义词词典的人力成本高，以英文为例，
据说现有的英文单词总数超过 1000 万个，光靠人力创造完全的同义词词典不太现实。
第三点是它无法表示单词的微妙差异，含义相近的单词之间可能也会有细微的区别，这些区别同义词词典不能很好体现。

\subsection{基于统计的方法}
基于统计的方法的目标是将单词表示成向量形式，在NLP中也称为分布式表示，这个想法基于分布式假设（distributional hypothesis）：
“某个单词的含义由它周围的单词形成”，也就是说单词含义可以它所在的上下文表示。
该方法需要使用语料库（corpus），即大量的文本数据，如Wikipedia等。
考虑句子："I say yes and you say no."，首先对文本进行分词，记录下每个单词的上下文（假设窗口大小为1），
I的上下文只有say，其表示为：
\begin{table}[H]
	\centering
	\begin{tabular}{c|cccccccccc}
		  & I & say & yes & and & you & no & . \\
		\hline
		I & 0 & 1   & 0   & 0   & 0   & 0  & 0 \\
	\end{tabular}
\end{table}
即you向量表示为$(0, 1, 0, 0, 0, 0, 0)$,继续处理剩余单词。最后生成一个生成共现矩阵（co-occurence matrix），
矩阵每行对应一个单词向量：

$$
	\begin{bmatrix}
		0 & 1 & 0 & 0 & 0 & 0 & 0\\
		1 & 0 & 1 & 0 & 1 & 1 & 0\\
		0 & 1 & 0 & 1 & 0 & 0 & 0\\
		0 & 0 & 1 & 0 & 1 & 0 & 0\\
		0 & 1 & 0 & 1 & 0 & 0 & 0\\
		0 & 1 & 0 & 0 & 0 & 0 & 1\\
        0 & 0 & 0 & 0 & 0 & 1 & 0
	\end{bmatrix}
$$

在上面的矩阵中，元素表示两个单词同时出现的次数，这种表示有时可能出现问题。
考虑语料库频繁出现“...the car...”的情况，如果只看单词的出现次数，
有可能出现 the 和 car 的相关性比 car 与 drive 相关性更强的情况。
为了解决这一问题，可以使用点互信息（Pointwise Mutual Information，PMI）指标，
对于随机变量 x 和 y，它们的 PMI 定义如下：
$$
\operatorname{PMI}(x, y) = \log \frac{P(x, y)}{P(x) P(y)}
$$
为了避免$\log 0$,实践上会使用正的点互信息（Positive PMI，PPMI）:
$$
\operatorname{PPMI}(x, y) = max (0, \operatorname{PMI}(x, y))
$$

\subsection{基于推理的方法}
基于统计的方法在处理大规模语料库会出现问题，由于单词数量非常大，生成的矩阵也会非常庞大，处理起来不太方便。
本节介绍基于推理的方法，他的原理相当于训练一个嵌入模型，让它做完形填空，通过反复的预测训练，学习到单词的嵌入表示。

在使用神经网络处理单词之前，首先需要将单词转换为向量表示。其中最简单的一种方式是使用 one-hot 编码。
这种表示方法中，向量的长度等于词汇表的大小，且仅将对应单词 ID 的位置设为 1，其余位置均为 0。

Word2Vec 是Google于2013开源的词向量计算工具，其中使用的模型之一为 CBOW（continuous bag-of-words）模型。
CBOW 模型的输入是上下文，这个上下文用单词列表表示，每个单词用one-hot表示成向量，因此如果上下文考虑 N 个单词，
则输入层会有 N 个。CBOW中，中间层的神经元是各个输入层经全连接层变换后得到的值的平均，输出层神经元数目等于词表大小，
这些神经元是各个单词的得分，它的值越大，说明对应单词的出现概率就越高。

Word2Vec 中还有一个是被称为 skip-gram 的模型，skip-gram 反转了 CBOW 模型的目标，CBOW模型从上下文预测单词，
而skip-gram模型则从单词预测上下文。在Word2Vec中，模型的输入侧的中间层权重就是单词的分布式表示。

考虑one-hot表示，由于其中只有一个为1的元素，它与中间层权重矩阵的乘积其实就是取矩阵的某一行，也即获得单词的分布式表示，
这个分布式表示也被称为词嵌入（word embedding），所以输入侧中间层有时也被称为 Embedding 层。由于这个操作比较简单，
实践上一般用查表的方式代替代替矩阵乘积。